\providecommand{\SO}{\mathrm{SO}}
\providecommand{\SE}{\mathrm{SE}}
\providecommand{\Eucl}{\mathrm{E}}
\providecommand{\Hyp}{\mathbb{H}}

\chapter{Applications}
\label{ch:applications}

The previous chapters built a theory: domains carry symmetries, symmetries
constrain architectures, and a single principle --- equivariant message
passing --- unifies convolutional networks, graph networks, and attention
mechanisms. This chapter turns to practice. Its purpose is not to catalogue
every system that uses geometric deep learning, but to show how the abstract
machinery of the ``5~Gs'' becomes concrete in real problems, and why the
geometric viewpoint translates into measurable gains in accuracy, data
efficiency, and physical consistency.

We organize the survey around three families of data. \emph{Computer vision}
ranges from regular pixel grids to unstructured three-dimensional point clouds,
exercising the grid, group, and gauge themes. \emph{Molecules and physical
systems} are the showcase of $\Eucl(3)$ and $\SE(3)$ equivariance, where
respecting the symmetries of physics is not a refinement but a requirement.
\emph{Recommender systems and social networks} are relational by nature and
rely on permutation invariance and, increasingly, on non-Euclidean geometry. A
fourth section, on \emph{geometric attention}, acts as a bridge: it shows how
the transformer --- the architecture behind much of modern machine learning ---
fits the geometric blueprint and reappears, suitably constrained, throughout the
other application areas.

\section{Computer vision}
\label{sec:apps:vision}

Computer vision is the historical birthplace of the ideas in this book.
Convolutional neural networks won the ImageNet competition in
2012~\citep{krizhevsky2012alexnet} precisely because they encode a symmetry of
images: a translation of the input produces a translated feature map. As
developed in \Cref{ch:grids}, this translation equivariance is what lets a
single learned filter be reused at every position, slashing the parameter count
and providing the inductive bias that makes images learnable from finite data.
Everything that follows in computer vision can be read as an attempt to extend
this success beyond the flat pixel grid.

\paragraph{Beyond translation.}
Many visual tasks are invariant to more than translation. A galaxy, a cell
under a microscope, or an aerial image has no preferred orientation; a
classifier should not have to relearn the same pattern in every rotated copy.
Group-equivariant convolutional networks (\Cref{ch:groups}) build rotation and
reflection symmetry directly into the convolution, so that a feature detected in
one orientation is automatically detected in all of them. The payoff is sharper
on small datasets, where data augmentation cannot cover the symmetry group
densely enough and the architecture must supply the invariance instead.

\paragraph{Three-dimensional vision.}
The cleanest break from the grid comes with three-dimensional data. A LiDAR
scan or a depth sensor returns a \emph{point cloud}: an unordered set of points
in $\R^3$, with no grid, no canonical ordering, and a variable cardinality. The
relevant symmetry is twofold. First, the output must be \emph{invariant to
permutations} of the points, since the set $\{p_1,\dots,p_n\}$ does not depend on
how we list its elements. Second, for many tasks it must be invariant or
equivariant to rigid motions of the scene. The first requirement is captured
exactly by the theory of functions on sets.

\begin{theorem}[Universal representation of set functions]
\label{thm:apps:deepsets}
Any continuous permutation-invariant function $f$ of a set
$\{x_1,\dots,x_n\}\subset\R^d$ can be written in the form
\[
  f(\{x_1,\dots,x_n\}) \;=\; \rho\!\left(\sum_{i=1}^n \phi(x_i)\right),
\]
for suitable continuous functions $\phi$ and $\rho$.
\end{theorem}

This result, due to the DeepSets framework~\citep{zaheer2017deepsets}, is the
theoretical foundation of architectures such as
PointNet~\citep{qi2017pointnet}: encode each point independently with a shared
network $\phi$, pool the encodings with a permutation-invariant aggregator (a
sum or a max), and decode with $\rho$. The construction is the simplest possible
instance of the message-passing blueprint of \Cref{ch:graphs} --- a single
round of messages on a graph with no edges --- and it already guarantees the
correct invariance by design rather than by augmentation.

\paragraph{Curved domains.}
When the data lives on a curved surface rather than in free space, the gauge and
geodesic themes (\Cref{ch:geodesics,ch:gauges}) take over. Spherical signals ---
panoramic images, cosmological maps, the surface of the Earth --- are processed
by spherical convolutional networks~\citep{cohen2018sphericalcnns} that are
equivariant to the rotation group $\SO(3)$, so that rotating the sphere rotates
the response coherently. On a general mesh there is no global frame, and
gauge-equivariant networks ensure that the output does not depend on the
arbitrary local choice of coordinates. In every case the recurring lesson is the
same: identify the symmetry of the domain, then build it into the operator.

\section{Geometric attention and transformers}
\label{sec:apps:attention}

The transformer~\citep{vaswani2017attentionneed} is, at first sight, a creature
of a different world: it was designed for sequences of words, not for grids or
graphs. Yet it fits the geometric picture precisely, and understanding why
explains both its success and the long line of geometric variants that power the
applications in the rest of this chapter.

\paragraph{Attention as adaptive message passing.}
Scaled dot-product attention computes, for each element, a weighted average of
the others,
\[
  \mathrm{Attention}(Q,K,V)
  = \operatorname{softmax}\!\left(\frac{QK^\top}{\sqrt{d_k}}\right) V,
  \qquad
  y_i = \sum_{j} \alpha_{ij}\, v_j,
\]
where the weights $\alpha_{ij}\ge 0$ sum to one and form a probability
distribution over the inputs. Read geometrically, this is message passing
(\Cref{ch:graphs}) in which the edge weights are not fixed by the domain but
\emph{computed from the data}: each element queries the others and decides how
much to listen to each. Where a convolution hard-codes the neighbourhood
weights and a graph network fixes them by the topology, attention learns them
on the fly. This places attention at the adaptive end of a spectrum whose rigid
end is the equivariant convolution --- a tension between built-in structure and
learned flexibility that recurs throughout geometric deep learning.

\paragraph{Attention on graphs.}
The most direct geometric instance is the Graph Attention Network
(GAT)~\citep{velickovic2017graph}, which restricts attention to the neighbours
of each node and weights them adaptively. For a node $i$ with neighbourhood
$\mathcal{N}(i)$, it computes attention coefficients from the node features,
normalizes them over the neighbourhood, and aggregates,
\[
  e_{ij} = \operatorname{LeakyReLU}\!\bigl(a^\top[Wh_i \,\|\, Wh_j]\bigr),
  \qquad
  \alpha_{ij} = \operatorname{softmax}_j(e_{ij}),
  \qquad
  h_i' = \sigma\!\Bigl(\textstyle\sum_{j\in\mathcal{N}(i)} \alpha_{ij}\,Wh_j\Bigr).
\]
Because the coefficients depend only on the features of the nodes involved and
not on any ordering, the layer is permutation-equivariant by construction
(\Cref{fig:apps:gat}). Like all standard message-passing networks, its power to
distinguish graphs is bounded by the Weisfeiler--Leman test~\citep{xu2019how},
but the learned, anisotropic weighting of neighbours makes it markedly more
expressive in practice than fixed-weight aggregation.

\begin{figure}[h]
\centering
\begin{tikzpicture}[>=Stealth, font=\small]
\tikzset{
    targetnode/.style={circle, draw=gdlred!70, fill=gdlred!20,
        minimum size=28pt, inner sep=0pt, font=\small\bfseries, line width=1.6pt},
    neighnode/.style={circle, draw=gdlblue!60, fill=gdlblue!20,
        minimum size=22pt, inner sep=0pt, font=\small},
    nonneighnode/.style={circle, draw=gdlgray!40, fill=gdlgray!15,
        minimum size=20pt, inner sep=0pt, font=\small, text=gdlgray!70!black}
}
\node[targetnode] (vi) at (0, 0) {$v_i$};
\node[neighnode] (v1) at (-2.4, 1.8)  {$v_1$};
\node[neighnode] (v2) at (0.0, 2.6)   {$v_2$};
\node[neighnode] (v3) at (2.4, 1.4)   {$v_3$};
\node[neighnode] (v4) at (-2.2, -1.4) {$v_4$};
\node[nonneighnode] (v5) at (2.6, -1.6)  {$v_5$};
\node[nonneighnode] (v6) at (-3.6, 0.2)  {$v_6$};
\draw[gdlgray!40, thin] (v1) -- (v6);
\draw[gdlgray!40, thin] (v4) -- (v6);
\draw[gdlgray!40, thin] (v3) -- (v5);
\draw[gdlgray!40, thin] (v1) -- (v2);
\draw[gdlgray!40, thin] (vi) -- (v1);
\draw[gdlgray!40, thin] (vi) -- (v2);
\draw[gdlgray!40, thin] (vi) -- (v3);
\draw[gdlgray!40, thin] (vi) -- (v4);
\draw[gdlgray!40, thin] (v5) -- (v3);
\draw[-{Stealth[length=7pt,width=5pt]}, line width=3.2pt,
    gdlorange, opacity=0.90, shorten >=8pt, shorten <=7pt] (v2) -- (vi);
\node[font=\scriptsize, text=gdlorange!70!black, fill=white, inner sep=1.5pt,
    rounded corners=1pt] at ($(v2)!0.48!(vi) + (0.7, 0.1)$) {$\alpha_{2i} = 0.42$};
\draw[-{Stealth[length=6pt,width=4.5pt]}, line width=2.2pt,
    gdlorange, opacity=0.75, shorten >=8pt, shorten <=7pt] (v1) -- (vi);
\node[font=\scriptsize, text=gdlorange!70!black, fill=white, inner sep=1.5pt,
    rounded corners=1pt] at ($(v1)!0.48!(vi) + (-0.85, -0.15)$) {$\alpha_{1i} = 0.28$};
\draw[-{Stealth[length=5.5pt,width=4pt]}, line width=1.5pt,
    gdlorange, opacity=0.60, shorten >=8pt, shorten <=7pt] (v3) -- (vi);
\node[font=\scriptsize, text=gdlorange!70!black, fill=white, inner sep=1.5pt,
    rounded corners=1pt] at ($(v3)!0.48!(vi) + (0.85, -0.15)$) {$\alpha_{3i} = 0.19$};
\draw[-{Stealth[length=5pt,width=3.5pt]}, line width=0.9pt,
    gdlorange, opacity=0.45, shorten >=8pt, shorten <=7pt] (v4) -- (vi);
\node[font=\scriptsize, text=gdlorange!70!black, fill=white, inner sep=1.5pt,
    rounded corners=1pt] at ($(v4)!0.48!(vi) + (-0.85, 0.15)$) {$\alpha_{4i} = 0.11$};
\begin{scope}[shift={(-3.6, -2.8)}]
    \node[font=\scriptsize\bfseries, text=gdlgray!70!black, anchor=west]
        at (0, 0) {Thickness $\propto$ attention weight};
    \draw[-{Stealth[length=4pt,width=3pt]}, line width=0.9pt,
        gdlorange, opacity=0.45] (0, -0.4) -- (1.2, -0.4);
    \node[font=\tiny, text=gdlgray!70!black, anchor=west] at (1.35, -0.4) {low};
    \draw[-{Stealth[length=5pt,width=4pt]}, line width=3.2pt,
        gdlorange, opacity=0.90] (2.4, -0.4) -- (3.6, -0.4);
    \node[font=\tiny, text=gdlgray!70!black, anchor=west] at (3.75, -0.4) {high};
\end{scope}
\end{tikzpicture}
\caption[Graph attention mechanism]{The graph attention mechanism. The target
node (red) computes attention coefficients $\alpha_{ij}$ with each neighbour
(blue); the arrow thickness reflects the magnitude of the learned weight. The
network learns to assign different importance to each neighbour from the node
features alone, so the operation is invariant to how the neighbours are
ordered.}
\label{fig:apps:gat}
\end{figure}

\paragraph{Equivariant attention in 3D.}
When the inputs carry geometric coordinates, attention must respect the
symmetries of space. The general principle is clean: an attention mechanism is
$G$-equivariant precisely when the compatibility scores are $G$-invariant and
the value transformation commutes with the group action.

\begin{proposition}[Equivariant attention]
\label{prop:apps:equiv-attn}
An attention layer is $G$-equivariant if the attention scores $\alpha_{ij}$ are
invariant under the action of $G$ and the value map commutes with that action.
Then $f(g\cdot x) = g\cdot f(x)$ for every $g\in G$.
\end{proposition}

The $\SE(3)$-Transformer~\citep{fuchs2020se3transformers} realizes this for
rigid motions in three dimensions. It computes attention scores from quantities
that are invariant under rotations and translations --- scalar features,
interatomic distances $\|r_i - r_j\|$, and inner products of vector features ---
and aggregates type-$\ell$ features so that vectors rotate with the scene while
scalars stay fixed. The Equiformer~\citep{liao2022equiformer} replaces the dot
product with an equivariant network acting on tensor products of irreducible
representations, allowing richer interactions while preserving $\SO(3)$
equivariance, and the Point Transformer~\citep{zhao2021point} brings vector
attention to point clouds. These models are the workhorses of the molecular and
physical applications of the next section.

\paragraph{Non-Euclidean and algebraic attention.}
Two further directions widen the reach of attention. Hyperbolic attention
networks~\citep{gulcehre2018hyperbolic} replace the Euclidean inner product with
operations in hyperbolic space, whose exponentially growing volume embeds
tree-like and hierarchical data with far less distortion than Euclidean space
(developed in \Cref{ch:geodesics}). Geometric Algebra Transformers
(GATr)~\citep{brehmer2023geometric} represent points, lines, and planes uniformly
as multivectors of a Clifford algebra and build $\Eucl(3)$-equivariant attention
on top, yielding a single architecture that handles diverse 3D geometric
primitives. The common thread is that the inductive bias of the geometry --- be
it negative curvature or the algebra of rigid motions --- is carried into the
attention operator itself.

\section{Molecules and physical systems}
\label{sec:apps:molecules}

If there is one domain that vindicates geometric deep learning, it is the
physical and life sciences. Molecules \emph{are} graphs whose nodes are atoms
and whose edges are bonds, embedded in $\R^3$; the laws governing them are
manifestly invariant under rigid motions of space. A model that ignores these
symmetries wastes capacity relearning physics that could simply be assumed.

\paragraph{Molecular property prediction.}
Graph neural networks predict pharmacological and quantum-chemical properties
directly from the molecular graph, replacing the hand-crafted descriptors and
string encodings (such as SMILES) of classical cheminformatics, which discard
the three-dimensional information that governs biological activity. When the
target depends on the molecular conformation --- energies, forces, dipole
moments --- the model must be $\Eucl(3)$- or $\SE(3)$-equivariant: rotating the
molecule must rotate predicted vectors and leave predicted scalars unchanged.
The equivariant attention architectures of the previous section, together with
tensor-field and equivariant message-passing networks, are precisely the tools
that enforce this.

\begin{example}[AlphaFold and protein structure]
\label{ex:apps:alphafold}
The most celebrated success of the field is
AlphaFold~\citep{jumper2021alphafold}, which solved the decades-old problem of
predicting a protein's three-dimensional structure from its amino-acid sequence.
At its core lies an attention-based architecture that reasons about the geometry
of residues in space while respecting the $\SE(3)$ symmetry of rigid motions:
the predicted structure transforms consistently when the frame of reference is
rotated or translated. Geometric equivariance is not a stylistic choice here but
a structural requirement, and it is a large part of why the model generalizes so
well.
\end{example}

\paragraph{Physical simulation.}
The same principles extend from static structures to dynamics. Equivariant
networks learn to simulate physical systems --- fluids, deformable materials,
many-body dynamics --- by operating on irregular meshes or particle sets while
respecting the symmetries of the underlying differential equations. Because the
symmetry is built in, the predictions are physically consistent by construction
(a rotated input yields a correspondingly rotated forecast) and the models often
match or exceed classical numerical solvers at a fraction of the cost. A single
foundation model can even be pretrained across small molecules and proteins at
once: equivariance is an architectural property, independent of the data, so a
network trained to be equivariant on one domain remains equivariant when
transferred to another with the same symmetry~\citep{jiao2024ept}.

\paragraph{Materials science.}
Crystalline solids add a further layer of symmetry: the discrete translational
and point-group symmetries of the crystal lattice. Encoding these symmetries in
equivariant networks lets models predict mechanical, thermal, and electronic
properties of candidate compounds, accelerating the search for new materials. As
in the molecular case, the geometric prior turns scarce, expensive labels into a
tractable learning problem.

\section{Recommender systems and social networks}
\label{sec:apps:graphs}

Relational data --- who interacts with whom, who buys what --- is graph-shaped
by nature, and here the governing symmetry is permutation invariance: the
identity of a user does not depend on the order in which we list the accounts.
Graph neural networks (\Cref{ch:graphs}) exploit this directly, propagating
information along the edges of the interaction graph so that each node's
representation is shaped by its neighbourhood.

\paragraph{Recommendation and social analysis.}
Recommender systems model users and items as nodes of a bipartite graph and
learn embeddings by message passing over observed interactions, capturing
collaborative signals that flat feature tables miss. Social-network analysis
uses the same machinery for node classification (inferring user attributes),
link prediction (suggesting connections), and community detection. In all of
these the irregular topology --- nodes with wildly different numbers of
neighbours and no natural ordering --- is exactly what convolutional models
cannot handle and graph networks are built for.

\paragraph{Hierarchies and hyperbolic geometry.}
Many relational datasets are not merely irregular but \emph{hierarchical}:
taxonomies, organizational charts, knowledge graphs, and the follower structure
of social networks all branch like trees. The number of nodes at distance $r$
from a root grows exponentially in $r$, and Euclidean space simply does not have
enough room to embed such growth without crushing distances. Hyperbolic space
does: its volume also grows exponentially, so a tree embeds with almost no
distortion (\Cref{fig:apps:hyperbolic}). This is the geometric motivation behind
hyperbolic embeddings and the hyperbolic attention networks of
\Cref{sec:apps:attention}, which place knowledge-graph entities and hierarchical
users in negatively curved space.

\begin{figure}[h]
\centering
\begin{tikzpicture}[
    root node/.style={circle, draw=gdlpurple!80, fill=gdlpurple!70,
        minimum size=12pt, inner sep=0pt},
    depth1 node/.style={circle, draw=gdlblue!80, fill=gdlblue!60,
        minimum size=10pt, inner sep=0pt},
    depth2 node/.style={circle, draw=gdlgreen!80, fill=gdlgreen!60,
        minimum size=10pt, inner sep=0pt},
    leaf node/.style={circle, draw=gdlyellow!80, fill=gdlyellow!70,
        minimum size=8pt, inner sep=0pt},
    tree edge/.style={draw=gdlgray!50, line width=0.8pt},
    scale=0.82, every node/.style={transform shape}
]
% LEFT: Euclidean
\begin{scope}[shift={(0,0)}]
    \coordinate (r) at (0, 1.9);
    \coordinate (d1-0) at (-1.8, 0.4); \coordinate (d1-1) at ( 1.8, 0.4);
    \coordinate (d2-0) at (-2.7, -1.1); \coordinate (d2-1) at (-0.9, -1.1);
    \coordinate (d2-2) at ( 0.9, -1.1); \coordinate (d2-3) at ( 2.7, -1.1);
    \coordinate (d3-0) at (-3.15, -2.6); \coordinate (d3-1) at (-2.25, -2.6);
    \coordinate (d3-2) at (-1.35, -2.6); \coordinate (d3-3) at (-0.45, -2.6);
    \coordinate (d3-4) at ( 0.45, -2.6); \coordinate (d3-5) at ( 1.35, -2.6);
    \coordinate (d3-6) at ( 2.25, -2.6); \coordinate (d3-7) at ( 3.15, -2.6);
    \draw[dashed, gdlgray!50, line width=0.8pt] (0,-0.35) circle (3.6);
    \draw[tree edge] (r) -- (d1-0); \draw[tree edge] (r) -- (d1-1);
    \draw[tree edge] (d1-0) -- (d2-0); \draw[tree edge] (d1-0) -- (d2-1);
    \draw[tree edge] (d1-1) -- (d2-2); \draw[tree edge] (d1-1) -- (d2-3);
    \draw[tree edge] (d2-0) -- (d3-0); \draw[tree edge] (d2-0) -- (d3-1);
    \draw[tree edge] (d2-1) -- (d3-2); \draw[tree edge] (d2-1) -- (d3-3);
    \draw[tree edge] (d2-2) -- (d3-4); \draw[tree edge] (d2-2) -- (d3-5);
    \draw[tree edge] (d2-3) -- (d3-6); \draw[tree edge] (d2-3) -- (d3-7);
    \node[root node] at (r) {};
    \node[depth1 node] at (d1-0) {}; \node[depth1 node] at (d1-1) {};
    \foreach \i in {0,...,3} {\node[depth2 node] at (d2-\i) {};}
    \foreach \i in {0,...,7} {\node[leaf node] at (d3-\i) {};}
    \node[font=\footnotesize\itshape, text=gdlgray!70!black] at (0, -4.6)
        {Euclidean embedding};
\end{scope}
% RIGHT: Poincaré disk
\begin{scope}[shift={(9.0, -0.35)}]
    \pgfmathsetmacro{\R}{3.6}
    \foreach \rad in {0.30, 0.52, 0.70, 0.83, 0.92} {
        \draw[gdlgray!20, thin] (0,0) circle (\rad*\R);}
    \foreach \ang in {0,30,60,90,120,150,180,210,240,270,300,330} {
        \draw[gdlgray!20, thin] (0,0) -- ({\ang}:\R);}
    \draw[black, line width=1.2pt] (0,0) circle (\R);
    \pgfmathsetmacro{\rone}{0.6044*\R}
    \pgfmathsetmacro{\rtwo}{0.8854*\R}
    \pgfmathsetmacro{\rthree}{0.9505*\R}
    \coordinate (pr) at (0, 0);
    \coordinate (pd1-0) at ({90}:{\rone}); \coordinate (pd1-1) at ({270}:{\rone});
    \coordinate (pd2-0) at ({50}:{\rtwo}); \coordinate (pd2-1) at ({130}:{\rtwo});
    \coordinate (pd2-2) at ({230}:{\rtwo}); \coordinate (pd2-3) at ({310}:{\rtwo});
    \coordinate (pd3-0) at ({35}:{\rthree}); \coordinate (pd3-1) at ({65}:{\rthree});
    \coordinate (pd3-2) at ({115}:{\rthree}); \coordinate (pd3-3) at ({145}:{\rthree});
    \coordinate (pd3-4) at ({215}:{\rthree}); \coordinate (pd3-5) at ({245}:{\rthree});
    \coordinate (pd3-6) at ({295}:{\rthree}); \coordinate (pd3-7) at ({325}:{\rthree});
    \draw[tree edge] (pr) -- (pd1-0); \draw[tree edge] (pr) -- (pd1-1);
    \draw[tree edge] (pd1-0) -- (pd2-0); \draw[tree edge] (pd1-0) -- (pd2-1);
    \draw[tree edge] (pd1-1) -- (pd2-2); \draw[tree edge] (pd1-1) -- (pd2-3);
    \draw[tree edge] (pd2-0) -- (pd3-0); \draw[tree edge] (pd2-0) -- (pd3-1);
    \draw[tree edge] (pd2-1) -- (pd3-2); \draw[tree edge] (pd2-1) -- (pd3-3);
    \draw[tree edge] (pd2-2) -- (pd3-4); \draw[tree edge] (pd2-2) -- (pd3-5);
    \draw[tree edge] (pd2-3) -- (pd3-6); \draw[tree edge] (pd2-3) -- (pd3-7);
    \node[root node] at (pr) {};
    \node[depth1 node] at (pd1-0) {}; \node[depth1 node] at (pd1-1) {};
    \foreach \i in {0,...,3} {\node[depth2 node] at (pd2-\i) {};}
    \foreach \i in {0,...,7} {\node[leaf node] at (pd3-\i) {};}
    \node[font=\footnotesize\itshape, text=gdlgray!70!black] at (0, -4.25)
        {Poincar\'e-disk embedding};
\end{scope}
\end{tikzpicture}
\caption[Euclidean versus hyperbolic embedding of a tree]{Embedding a balanced
tree in Euclidean space (left) and in the Poincar\'e disk (right). In the plane,
successive levels crowd together and sibling subtrees overlap, distorting
distances; in hyperbolic space the exponentially growing volume gives every
level room to spread out, so the tree embeds with little distortion. This is why
hierarchical relational data is naturally modelled in negatively curved space.}
\label{fig:apps:hyperbolic}
\end{figure}

\paragraph{Temporal graphs: fraud and risk.}
Financial transactions form graphs whose topology evolves in time. Graph neural
networks operating on these temporal graphs detect fraud rings and systemic
risk by recognizing structural patterns --- unusual subgraphs, sudden changes in
connectivity --- that tabular models, blind to the relational structure, cannot
see. The geometric prior is again decisive: it is the \emph{shape} of the
interaction pattern, not the individual features, that carries the signal.

\section{A common principle}
\label{sec:apps:common}

Across these domains a single message recurs. In computer vision the symmetry is
translation, rotation, or the gauge freedom of a curved surface; in molecular
and physical modelling it is the Euclidean group of rigid motions; in relational
data it is permutation of nodes, and the geometry of the embedding space may be
flat or negatively curved. In every case the same recipe applies: identify the
symmetry the problem respects, build it into the architecture, and let the
network spend its capacity on what actually varies.

\begin{remark}[Geometry as the organizing principle]
\label{rem:apps:unification}
The breadth of these applications is itself the strongest evidence for the
central claim of this book. Convolutional networks for images, equivariant
transformers for proteins, and graph networks for social data look like
unrelated tools, yet each is an instance of equivariant message passing,
differing only in the domain and the symmetry group it respects. Geometry is not
a domain-specific trick but the common principle that geometric deep learning
identifies, formalizes, and exploits.
\end{remark}

Beyond these consolidated areas, the same ideas are opening lines of research
wherever geometric structure is fundamental: computational neuroscience, where
graph networks model structural and functional brain connectivity and reveal
patterns of neurological disorder invisible to tabular methods; and climate
science, where the spherical geometry of the Earth invites $\SO(3)$-equivariant
spherical networks for weather and climate prediction. The diversity of these
settings underscores that respecting the geometry of data is not an
optimization detail but the organizing principle of the field, and the bridge
from the theory of the preceding chapters to the systems that are reshaping
science and technology.
